% Appendix B: Digital Resources and Computational Toolchain

This monograph is accompanied by an open suite of machine-readable datasets, relational polygon manifests, and interactive command-line utilities. All digital artifacts are designed to facilitate computational verification, algorithmic querying, and combinatorial research into cyclic lattice composition mosaics.

\section{The Atlas Index}
\label{sec:atlas-index}

The complete inventory of all $21{,}030$ affine macro-orbits is compiled in the machine-readable dataset:
\begin{center}
\path{unimosatlas/data/hexadecality_index.txt}
\end{center}
This comprehensive index contains $21{,}065$ lines covering the entire combinatorial universe of $83{,}488$ dihedral orbits ($14$ Monomers, $295$ Dimers, and $20{,}721$ Tetramers):

\begin{enumerate}
    \item \textbf{Record Identifier}:
    Prefixed by orbit class and index: \texttt{M01}--\texttt{M14} for monomers, \texttt{D001}--\texttt{D295} for dimers, and \texttt{T00001}--\texttt{T20721} for tetramers.
    
    \item \textbf{Physical Metric Invariants}:
    Each record tabulates the dihedral orbit size $|O| \in \{1, 4, 16, 32\}$, the block bilateral symmetry count $s \in \{0, \dots, 4\}$, and the total chord energy $E$.
    
    \item \textbf{Partition-Slash Notation}:
    Partitions are rendered in compact hexadecimal slash notation, where vertices $10 \dots 15$ are represented by uppercase letters $\text{A} \dots \text{F}$. For example:
    \begin{verbatim}
    /0279/18AF/35CE/46BD/
    \end{verbatim}
    denotes the partition $\{\{0, 2, 7, 9\}, \{1, 8, 10, 15\}, \{3, 5, 12, 14\}, \{4, 6, 11, 13\}\}$.
    
    \item \textbf{Block Chirality Signatures}:
    Appended to each partition is a 4-letter signature in brackets, such as \texttt{[ssss]}, \texttt{[ssnn]}, or \texttt{[nnnn]}. Each character corresponds to one block:
    \begin{itemize}
        \item \texttt{s} (symmetric): the block possesses an internal bilateral reflection axis in $\mathbb{Z}_{16}$, rendering it achiral and non-orientable (shaded blue in plates).
        \item \texttt{n} (non-symmetric): the block lacks any internal reflection axis, rendering it intrinsically chiral (shaded pink in plates).
    \end{itemize}
    
    \item \textbf{Galois Dilation Cycle Transitions}:
    For dimers and tetramers, the line tracks the transition cycle under the affine generator $\mu_3 \colon x \mapsto 3x \pmod{16}$:
    \begin{verbatim}
    D00001 Dimer |O|=4 s=4 E=640 /028A/139B/46CE/57DF/ [ssss]
      -> /068E/179F/24AC/35BD/ [ssss]
    \end{verbatim}
\end{enumerate}

\section{The Mosaical Orbit Resolver}
\label{sec:lookup-tool}

The command-line tool \texttt{unimosatlas/tools/lookup\_mosaic.py} provides an automated resolution engine for any user-supplied set partition of $\mathbb{Z}_{16}$. 

Given an arbitrary partition string in partition-slash format, the tool:
\begin{enumerate}
    \item Verifies that the string forms a valid, non-overlapping partition of $\{0, \dots, 15\}$ into four blocks of size four.
    \item Computes the full $32$-element dihedral orbit under $D_{16}$ and extracts the lexicographical canonical representative.
    \item Identifies the stabiliser subgroup and determines the dihedral orbit size $|O| \in \{1, 4, 16, 32\}$.
    \item Calculates the chord energy $E$ and the block chirality signature.
    \item Resolves the parent macro-orbit (Monomer, Dimer, or Tetramer) and its primary catalogue index.
    \item Traces the full macro-orbit under Galois dilation $\mu_3 \colon x \mapsto 3x \pmod{16}$, printing both the geometrically dilated chord diagrams and the corresponding canonical dihedral indices.
\end{enumerate}

\begin{example}
Querying the tetramer \texttt{/0279/18AF/35CE/46BD/} yields:
{\small
\begin{verbatim}
$ lookup-mosaic "/0279/18AF/35CE/46BD/"
======================================================================
HEXADECALITY: COMPOSITION MOSAIC ORBIT & CARTOUCHE RESOLVER
======================================================================
Input Partition:       /0279/18AF/35CE/46BD/  [ssss]
Canonical Dihedral D16:/0279/18AF/35CE/46BD/  [ssss]
Dihedral Orbit Size:   |O| = 4 (Stabiliser order: 8)
Symmetric Blocks:      s = 4 of 4 (Achiral)
Chord Energy:          E = 720
----------------------------------------------------------------------
Affine Macro-Orbit Type: Tetramer (4-Cycle)

Orbital Comosaics Cycle (Physical Dilation vs Canonical Rep):
  Phase 1 (mu_1 = id) : /0279/18AF/35CE/46BD/ [ssss] (Canonical)
  Phase 2 (mu_3 = x3) : /056B/127C/38DE/49AF/ [ssss]
                        (Canonical: /016B/27CD/389E/45AF/)
  Phase 3 (mu_7 = x7) : /01EF/2345/6789/ABCD/ [ssss]
                        (Canonical: /0123/4567/89AB/CDEF/)
  Phase 4 (mu_5 = x5) : /03AD/147E/258B/69CF/ [ssss]
                        (Canonical: /0369/18BE/25CF/47AD/)
======================================================================
\end{verbatim}}
\end{example}

\section{The Relational Polygon Datasets and Embedding Manifests}
\label{sec:hexdec-relational-datasets}

To establish a bidirectional bridge between composition mosaics and the canonical inscribed polygon census of \texttt{polyatlas}, the entire combinatorial corpus across all three metric universes ($\mathbb{Z}_4$, $\mathbb{Z}_9$, and $\mathbb{Z}_{16}$) is provided as normalised, machine-readable tab-separated resource files housed in \path{unimosatlas/data/}.

\subsection{Relational Architecture and Syntax}

Each universe is structured as a relational pair:
\begin{enumerate}
    \item \textbf{Mosaic-Level Catalogue (\path{data/z*_mosaics.tsv})}:
    Tabulates the global dihedral and affine invariants ($|O|$, $|\operatorname{Stab}|$, $s$, $E$, chirality signature), the constituent polygon multiset profile, and the collective \emph{block manifest}.
    
    \item \textbf{Block-Level Relational Bridge (\path{data/z*_mosaic_blocks.tsv})}:
    Links each individual block in every canonical mosaic directly to its parent inscribed polygon class in \path{polyatlas/data/inngons.tsv}.
\end{enumerate}

Every block embedding is encoded by the standardised manifest descriptor:
\begin{equation}
\label{eqn:hexdec-block-manifest}
\text{[Class]}@\text{[Root]}[\text{Chirality}]:\text{[Realised Strides]}
\end{equation}
where \textbf{Class} denotes the parent polygon class $k\text{-}i$ defined in \texttt{polyatlas}; \textbf{Root} is the starting vertex $v_0 = \min(B_j)$ around the clock face; \textbf{Chirality} specifies the orientation sign ($+$ for direct rotational alignment, $-$ for retrograde reflection); and \textbf{Realised Strides} lists the forward perimeter strides from $v_0$ along the block.

\subsection{Multi-Universe Relational Directory}

The companion datasets span three cyclic universes:
\begin{itemize}
    \item \textbf{$D_4$ Universe ($z_2^2$ on $\mathbb{Z}_4$)}:
    \path{data/z4_mosaics.tsv} ($2$ orbits) and \path{data/z4_mosaic_blocks.tsv} ($4$ block records). Both digon classes ($2\text{-}1$ and $2\text{-}2$) are realised.
    
    \item \textbf{$D_9$ Universe ($z_3^3$ on $\mathbb{Z}_9$)}:
    \path{data/z9_mosaics.tsv} ($25$ dihedral comosaics across $9$ macro-orbits) and \path{data/z9_mosaic_blocks.tsv} ($75$ block records). All $7$ triangle classes ($3\text{-}1$ through $3\text{-}7$) are fully realised.
    
    \item \textbf{$D_{16}$ Universe ($z_4^4$ on $\mathbb{Z}_{16}$)}:
    \begin{itemize}
        \item \textbf{Primary Macro-Orbits}:
        \path{data/z16_macro_mosaics.tsv} ($21{,}030$ records) and\\
        \path{data/z16_macro_mosaic_blocks.tsv} ($84{,}120$ block records).
        \item \textbf{Exhaustive Dihedral Comosaics}:
        \path{data/z16_mosaics.tsv} ($83{,}488$ records) and\\
        \path{data/z16_mosaic_blocks.tsv} ($333{,}952$ block records).
    \end{itemize}
\end{itemize}

\subsection{The Injection Property and Cascade Filter}

Unlike unconstrained partitions, the mosaics treated here are equal-block partitions ($z_k^m$), representing structured injections into the general compositional universe:
\begin{itemize}
    \item \textbf{Enantiomorphic Realisation}: In \texttt{polyatlas}, chiral polygons appear strictly in their canonical direct orientation. In mosaics, chiral blocks appear in both direct ($+$) and retrograde ($-$) configurations (e.g.\ class $4\text{-}68$ exhibits $3{,}303$ direct and $2{,}563$ retrograde blocks across $\mathbb{Z}_{16}$).
    \item \textbf{Universal Surjectivity}: Across the complete $83{,}488$ comosaics of $\mathbb{Z}_{16}$, all $72$ tetragon classes from \texttt{polyatlas} are realised ($100\%$ coverage).
    \item \textbf{Sorting Cascade Dispersion}: In the $21{,}030$ primary macro-orbit representatives, exactly $70$ of the $72$ classes appear. Classes $4\text{-}1$ ($/0123/$, strides $111\text{D}$) and $4\text{-}3$ ($/0134/$, strides $121\text{C}$) are filtered out of the leading phase because the 5-tier natural sorting cascade prioritises high chord dispersion energy $E$, systematically shifting clustered blocks into dilated comosaic phases ($\mu_3$, $\mu_7$, or $\mu_5$).
\end{itemize}

\begin{table}[htbp]
\centering
\caption{\label{tab:hexdec-sample-manifests}Exemplar Composition Mosaics and their Relational Manifests}
\resizebox{\linewidth}{!}{%
\begin{tabular}{l c c c l l}
\toprule
Identifier & Type & $|O|$ & $s$ & Polygon Profile & Relational Embedding Manifest \\
\midrule
$M_1$ ($\mathbb{Z}_4$) & Monomer & 1 & 2 & $2 \times 2\text{-}2$ & \path{2-2@0+:22/2-2@1+:22} \\
$M_2$ ($\mathbb{Z}_4$) & Monomer & 2 & 2 & $2 \times 2\text{-}1$ & \path{2-1@0+:13/2-1@2+:13} \\
\midrule
$M_1$ ($\mathbb{Z}_9$) & Monomer & 1 & 3 & $3 \times 3\text{-}4$ & \path{3-4@0+:333/3-4@1+:333/3-4@2+:333} \\
$T_1$ ($\mathbb{Z}_9$) & Trimer & 3 & 3 & $3 \times 3\text{-}2$ & \path{3-2@0+:144/3-2@2+:414/3-2@3+:144} \\
$T_3$ ($\mathbb{Z}_9$) & Trimer & 9 & 1 & $3\text{-}2 + 2 \times 3\text{-}7$ & \path{3-2@0+:144/3-7@2+:234/3-7@3-:324} \\
\midrule
$M_{01}$ ($\mathbb{Z}_{16}$) & Monomer & 1 & 4 & $4 \times 4\text{-}28$ & \path{4-28@0+:4444/4-28@1+:4444/4-28@2+:4444/4-28@3+:4444} \\
$M_{02}$ ($\mathbb{Z}_{16}$) & Monomer & 4 & 4 & $2 \times 4\text{-}13 + 2 \times 4\text{-}27$ & \path{4-13@0+:1717/4-27@2+:5353/4-27@3+:3535/4-13@4+:1717} \\
$M_{05}$ ($\mathbb{Z}_{16}$) & Monomer & 16 & 4 & $4\text{-}13+4\text{-}22+4\text{-}27+4\text{-}28$ & \path{4-13@0+:1717/4-27@2+:3535/4-28@3+:4444/4-22@4+:2626} \\
$T_{00001}$ ($\mathbb{Z}_{16}$) & Tetramer & 4 & 4 & $4 \times 4\text{-}12$ & \path{4-12@0+:1636/4-12@2+:1636/4-12@3+:3616/4-12@4+:1636} \\
\bottomrule
\end{tabular}%
}
\end{table}
